\documentclass{letter}
\usepackage[T1]{fontenc}
\usepackage{lmodern}

\signature{Maximiliano Martino and Daniel Fraiman}
\address{Universidad de San Andr\'es\\Argentina}

\begin{document}
\begin{letter}{Editors\\Journal of Statistical Software}

\opening{Dear Editors,}

Thank you for granting permission to resubmit a substantially revised version
of our manuscript. We took the editorial assessment as a request for a
software and inferential redesign, rather than a request for cosmetic changes.

The revised manuscript is entitled ``BrainNetTest: Global and Edge-Wise
Inference for Binary Brain Network Populations in R.'' The package now uses
validated S3 input and result classes, standard print, summary, plot, and
data-frame methods, and a deliberately smaller public API. Global
randomization inference is separated from multiplicity-controlled edge
inference. The prior adaptively stopped ``critical edge'' procedure is no
longer presented as inferential; its retained prefix-sum path is explicitly
descriptive and its computational gain is benchmarked.

We also rebuilt the replication materials. The standalone
\texttt{code-full.R} script regenerates every reported simulation and
benchmark result, table, figure, and numeric macro, while a separate short
\texttt{code.R} and \texttt{code.html}
exercise a reduced study. Both scripts use only the installed package and
relative paths. The included README maps all outputs, and the full run records
session and provenance information.

The manuscript is now organized around the clinical brain-network problem,
statistical method, software contract, classes, workflows, real-data
application, validation, and computational performance. The ABIDE case study
compares autism and control resting-state connectomes in both the full cohort
and a sex-, age-, and motion-balanced sensitivity sample. In both samples the
observed global statistic is extreme under label randomization, while no
individual connection survives strict multiplicity control. The manuscript
therefore reports a preprocessing-conditional population association rather
than a causal, diagnostic, or biomarker claim. The standalone application
script and derived data make this analysis reproducible without access to raw
images.

A detailed point-by-point response accompanies this letter.

\closing{Sincerely,}
\end{letter}
\end{document}
